\section{Experimental Results}

\begin{figure*}[ht]
    \centering
    \begin{overpic}
    [width=0.95\linewidth]{figures/comparison_general.pdf}
    \put(5, -1){Text prompts}
    \put(22,-1){(a) GraphDreamer}
    \put(44,-1){(b) Blender-MCP}
    \put(67,-1){(c) MIDI}
    \put(86,-1){(d) Ours}
    \end{overpic}
    \vspace{0.2cm}
    \caption{\textbf{Comparison to baseline methods.} The scenes are generated from our text prompts. OOM indicates out of memory.}
    \label{fig:comparison_general}
\end{figure*}





\subsection{Comparison}

\begin{table}[]
    \centering
    \begin{tabular}{c| ccccc }
        \toprule
            & Clip Score $\uparrow$ & VQA Score $\uparrow$ & Displacement $\downarrow$ & Pene. Ratio $\downarrow$ & Phys. Score $\uparrow$ \\
        \hline
        GraphDreamer & 27.53 & 0.46  & 0.25 & 61.72 & 40.0 \\
        Blender-MCP & 28.93 & 0.56 & 1.03 & 14.78 & 47.7 \\
        MIDI & 29.68 &  0.63  & 0.69 & 110.80 & 62.7  \\
        \hline
        Raw layout & 29.88 & 0.64  & 0.81 & 14.11 & 65.5 \\
      Scene Init. & 30.77 & \textbf{0.70} & 2.91 & \textbf{0} & 34.2 \\
        \hline
        Ours & \textbf{31.79} &  0.68  & \textbf{0} & \textbf{0} & \textbf{88.5}  \\
        \bottomrule
    \end{tabular}
    \caption{\textbf{Quantitative Evaluation.} Our method achieves the highest semantic consistency with input text prompts among all baselines, and is the only method that achieves perfect physical stability and non-intersection. We also provide ablation results: removing both layout initialization and optimization (shown as \textit{raw layout}), and removing only layout optimization (shown as \textit{scene init.}).
    }
    \label{tab:metrics}
    \vspace{-0.5cm}
\end{table}






\subsubsection{Baselines}

We compare our method against three baselines: GraphDreamer \cite{gao2024graphdreamer}, MIDI \cite{huang2024midi}, and Blender-MCP \cite{blender_mcp}. Both GraphDreamer and Blender-MCP take text prompts as input, while MIDI uses a reference image as input. To ensure a fair comparison, we provide MIDI with our scene reference image as their input. 



\subsubsection{Dataset}

Since there is no standard benchmark for general scene generation, we construct our own test dataset consisting of 18 text prompts. Among them, 3 prompts are taken from MIDI, and 2 prompts are from GraphDreamer. Additionally, we use an LLM to generate 13 new text prompts spanning diverse scenes. These prompts describe physical interactions between objects, including a stack of books and a basket of fruits. Additional 3D results generated by our method, in comparison with the baselines, along with corresponding text prompts and reference images, can be found on visualization website\footnote{\url{https://3dsim-baseline-visualization.netlify.app/}}. Beyond these comparisons, we also present 12 more examples produced by our method in \autoref{supp:more_examples}.

\subsubsection{Evaluation Metrics}

We evaluate our generated scenes using five metrics: CLIP Score\cite{radford2021learningtransferablevisualmodels}, VQA Score\cite{lin2024evaluatingtexttovisualgenerationimagetotext}, Simulated Scene Displacement (D), the Ratio of Penetrating Triangle Pairs (R), and a Physical Plausibility Score. Together, these metrics measure semantic consistency, physical stability, interpenetration, and overall physical plausibility. Details of all metrics are provided in \autoref{supp:metrics}.

\subsubsection{Performance and Discussions}

In \autoref{fig:comparison_general}, we compare PAT3D with baseline methods on five general scenes involving complex object interactions. Additional comparisons with four scenes highlighted in the MIDI and GraphDreamer are provided in \autoref{supp:more_examples}. Importantly, in PAT3D, reference image serves only to extract the complex spatial relations implied in the text prompt; the final scene does not need to remain visually consistent with the reference image. 

GraphDreamer struggles to scale to larger scenes because it jointly optimizes both object geometry and scene layout through Score Distillation Sampling (SDS)~\cite{poole2022dreamfusion}, which is highly resource-intensive. Moreover, GraphDreamer exhibits weak understanding of spatial relations in text prompts. As shown in the second, fourth, and fifth scenes of \autoref{fig:comparison_general}, it often ignores spatial constraints, leading to chaotic object arrangements. Blender-MCP generates layouts with little physical realism. In the first and second scenes of \autoref{fig:comparison_general}, the razor and toothbrush float above the cup, and the plate is suspended in mid-air. It also produces objects with unrealistic scales: in the fourth scene, the cake and vase appear disproportionately small compared to the table. MIDI encounters difficulties when handling scenes with complex object contact, as seen in the first, second, and fifth scenes of \autoref{fig:comparison_general}, objects often appear in irregular yet tightly packed configurations. Although interpenetration is avoided, the resulting layouts are cluttered potentially because MIDI generates the entire scene in a single step, the quality of individual objects is compromised.

By contrast, our method decomposes the scene generation process, iteratively creating objects to ensure high-quality results. Leveraging VLM-based guidance together with physics simulations, PAT3D produces 3D scenes that are both physically realistic and semantically consistent, even in scenarios with complex object interactions. Quantitative comparisons are shown in \autoref{tab:metrics}. Compared to baseline methods, which frequently suffer from object intersection and floating artifacts that undermine physical plausibility, our approach consistently produces stable, penetration-free arrangements. Our method also achieves the highest semantic consistency with the input text prompts.

\subsection{Application}

Our generated simulation-ready scenes can be directly imported into a simulator for downstream applications. We demonstrate two such applications: scene editing and robotic manipulation.

\subsubsection{Scene Editing}




We demonstrate a scene editing application enabled by our framework, which supports interactive manipulation while preserving the physical plausibility of the entire scene, including object addition and deletion. By leveraging our physics-based simulation backend, the edited scene converges to a force-equilibrium state without mesh intersections. \autoref{fig:animation} highlights an example showcasing object addition and deletion.

\begin{figure}[ht]
    \centering
         \begin{overpic}
     [width=0.95\linewidth]{figures/scene_editing.pdf}
     \put(10,-3){(a)}
     \put(35,-3){(b)}
     \put(60,-3){(c)}
     \put(85,-3){(d)}
     \end{overpic}
        \vspace{0.2cm}
        \caption{\textbf{Scene editing.} We demonstrate the equilibrium state after addition and deletion operations: (a) initial scene, (b) deleting a book at the bottom, (c) deleting the pen holder, (d) adding a book on top.}
    \label{fig:animation}
\end{figure}

\subsubsection{Robotic Manipulation}
\begin{wrapfigure}{r}{0.44\linewidth}
    \centering
    \vspace{-0.8cm}
    \begin{overpic}[width=\linewidth]{figures/robot_manipulate.pdf}
    \put(5,-5){Successful grasp}
    \put(55,-5){Failed grasp}
    \end{overpic}
    \caption{\textbf{Policy evaluation for robotic manipulation.} 
    Example of a successful and a failed grasp where the attempted action causes objects to topple.}
    \vspace{-0.5cm}
    \label{fig:robot_manipulate}
\end{wrapfigure}
Our generated scenes can be directly imported into a simulator to validate robotic manipulation policies. In \autoref{fig:robot_manipulate}, we present two illustrative examples, a failed grasp and a successful grasp. Robotic manipulation applications impose unique requirements on scene generation: objects must be consistently positioned and free of interpenetrations. Our framework satisfies these requirements, ensuring that the generated scenes are well-suited for reliable policy evaluation.







\subsection{Ablation Study}\label{sec:ablation}

\begin{figure*}[ht]
    \vspace{-0.5cm}
    \centering
    \begin{minipage}[t]{0.46\linewidth}
        \centering
        \begin{overpic}[width=\linewidth]{figures/ablation_scene_tree.pdf}
            \put(5,-3){(a) Depth prediction}
            \put(57,-3){(b) Scene tree}
        \end{overpic}
        \caption{\textbf{Layout initialization w/o and w/ scene tree.} Layouts obtained from depth prediction (a) without and (b) with adjustment based on the scene tree.
        (Text prompt: \textit{``...a neatly stacked pile of three books...''}. See \autoref{supp:text_prompt} for the complete prompt.)}\label{fig:ablation_scene_tree}
    \end{minipage}
    \hspace{0.8cm}
    \begin{minipage}[t]{0.46\linewidth}
        \centering
        \begin{overpic}[width=\linewidth]{figures/ablation_diff_sim.pdf}
            \put(15,-0.5){(a) Before}
            \put(67,-0.5){(b) After}
        \end{overpic}
        \caption{\textbf{Layout optimization.} Simulated layouts from initial layout (a) without and (b) with further optimization using differentiable simulation.
    (Text prompt: \textit{``a stack of colorful wooden blocks...''}. See \autoref{supp:text_prompt} for the complete prompt.)
    }\label{fig:ablation_diff_sim}
    \end{minipage}
\end{figure*}


We first qualitatively illustrate the impact of our layout initialization module and simulation-in-the-loop optimization module. As shown in \autoref{fig:ablation_scene_tree}(a), while the spatial relations between objects extracted directly from the depth map are generally reasonable, the scene still suffers from significant interpenetration: books intersecting with one another and the pen protruding its holder. 
By contrast, after applying our proposed layout initialization based on the scene tree, we obtain an intersection-free layout shown in \autoref{fig:ablation_scene_tree}(b), where projections of each object along gravity direction typically lies in the projections of their containers or supporters, thereby satisfying physical dependency constraints along the gravity direction.

Nevertheless, simply enforcing physical dependencies before simulation does not ensure that the resulting simulated scene would still satisfy the intended semantics. In \autoref{fig:ablation_diff_sim}(a), a stack of irregular blocks collapses after simulation due to an unbalanced center of mass. In contrast, \autoref{fig:ablation_diff_sim}(b) demonstrates that, by further optimizing the initial layout through our simulation-in-the-loop optimization, the simulated scene converges to a stable configuration of stacked blocks that better reflects the semantics.

We further quantitatively evaluate the effectiveness of our method by computing semantic and physical metrics on both the depth-aligned layout without our layout initialization and optimization (denoted as \textit{raw layout}), the layout from our scene initialization module (denoted as \textit{{Scene Init.}}), and compare them with our final output in \autoref{tab:metrics}. Our layout initialization removes penetration by sacrificing physical stability, but it enables applying our layout optimization to consistently improve all metrics.  The gains in semantic consistency metrics are relatively smaller, as they primarily depend on visual appearance factors such as geometry and texture.










